\documentclass{situationist-preprint}

\title{Internalism for Rational Agents}
\msid{MTE-07-a}
\msfamily{MTE-07}
\msdiscipline{Metaethics}
\mstype{Principal manuscript}
\mskeywords{motivational internalism; amoralist; practical rationality; moral judgment; reasons; acquired sociopathy}

\begin{document}
\maketitle

\begin{abstract}
Does moral judgment necessarily motivate? Internalists say that it does, so that an agent who sincerely judges that she ought to act and is wholly unmoved is impossible; externalists say that the connection is contingent, and point to the amoralist, who makes moral judgments and does not care. We defend a conditional internalism: necessarily, a practically rational agent who judges that she morally ought to act is motivated, to some degree, to act accordingly. Amoralists are possible, but they are practically irrational. The argument has two premises. The first is that to judge that one morally ought to act is to judge that there are decisive moral reasons for one to act, which follows from the account of ought in terms of reasons. The second is that a practically rational agent is moved, to some degree, by what she takes to be reasons for her. The amoralist's judgment is a belief, and beliefs and motivations can come apart; so the amoralist is possible. What she lacks is a response that her own belief calls for, and that is a failure of practical rationality.
\end{abstract}

\section{The problem}

Hume argued that morals excite passions and produce or prevent actions, that reason alone cannot do this, and that moral distinctions are therefore not derived from reason \citep[3.1.1]{HumeTreatise}. The argument assumes what has come to be called \emph{motivational internalism}: that there is a necessary connection between moral judgment and motivation. \citet{Frankena1958} gave the debate its modern shape by distinguishing internalist theories, on which motivation is built into the content or nature of moral judgment, from externalist ones, on which it is supplied by a desire that the judgment does not contain.

Internalism comes in strengths. On the strongest form, anyone who sincerely judges that she ought to act is motivated to act; this is the form expressivists have relied on, since it is explained if moral judgments are states of motivation rather than beliefs \citep{Hare1952}. Externalists reply with the amoralist, who understands moral language, judges that some acts are right and others wrong, and is unmoved. \citet{Brink1989} argued that the amoralist is conceivable, and that internalism must therefore be false or trivial; \citet{Svavarsdottir1999} argued that the evidence for internalism is better explained by the externalist hypothesis that most people have a standing desire to do what is right. The internalist's reply has been that the amoralist does not really make moral judgments but uses moral terms in inverted commas, as when one says that something is right meaning that it is what others call right \citep{Hare1952,Smith1994}.

The empirical literature has been brought to bear. Patients with damage to the ventromedial prefrontal cortex appear to make normal moral judgments while lacking the motivational and emotional responses that normally accompany them \citep{Roskies2003}, and psychopaths are often described as understanding what morality requires without caring \citep{Nichols2002}. If these agents make genuine moral judgments, strong internalism is false.

A weaker form of internalism survives these cases. \citet{Korsgaard1986} argued that practical reason motivates those who are rational, and that failures of motivation are failures of rationality. \citet{Smith1994} defended the practicality requirement: if an agent judges that it is right for her to act in certain circumstances, then either she is motivated to act or she is practically irrational. We defend a version of this view, with an argument that rests on the account of ought in terms of reasons.

\begin{principle}[Thesis]\label{p:thesis}
Necessarily, if a practically rational agent judges that she morally ought to $\phi$, she is motivated, to some degree, to $\phi$. There can be agents who judge that they morally ought to $\phi$ and are not motivated to $\phi$; such agents are practically irrational.
\end{principle}

\section{Framework}

\begin{definition}[Moral judgment]\label{d:judgment}
A \emph{moral judgment} that one ought to $\phi$ is the belief that one morally ought to $\phi$.
\end{definition}

The definition takes sides on cognitivism, as MTE-01 does: moral facts are natural facts, and judgments about them are beliefs. The internalism we defend does not depend on any view of the content of these beliefs beyond the one stated in \ref{p:reasons}.

\begin{definition}[Motivation]\label{d:motivation}
An agent is \emph{motivated to some degree} to $\phi$ if she has a disposition to $\phi$, or to intend to $\phi$, that would issue in action in the absence of stronger contrary motivation.
\end{definition}

Motivation in this sense is pro tanto and defeasible. An agent motivated to some degree to keep her promise may break it because she is more strongly moved to do something else; ACT-04 argues that weakness of will is possible, and the akratic agent who acts against her judgment is motivated by it, though not enough.

\begin{principle}[Ought and reasons]\label{p:reasons}
To morally ought to $\phi$ is for the moral reasons for $\phi$-ing to be stronger than the moral reasons for each alternative; and moral reasons are reasons.
\end{principle}

\ref{p:reasons} is the definition of ought in terms of reasons that MTE-11 defends, restricted to moral reasons. It has a consequence for the content of moral judgment. An agent who believes that she morally ought to $\phi$, and who possesses the concept of moral ought, believes something that entails that there are reasons for her to $\phi$; and if she grasps the entailment, as anyone who possesses the concept does, she believes that there are reasons for her to $\phi$. An agent who believes that she morally ought to $\phi$ and also believes that there is no reason at all for her to $\phi$ has inconsistent beliefs.

\begin{definition}[Practical rationality]\label{d:rational}
An agent is \emph{practically rational}, in the sense at issue, if her motivations respond to the considerations represented by her own beliefs about her reasons: if she believes that there is reason for her to $\phi$, she is motivated to some degree to $\phi$.
\end{definition}

Definition~\ref{d:rational} specifies one member of the family of senses of rationality that MTE-11 defines: rationality relative to the considerations an agent's beliefs represent, the sense in which, following \citet{Parfit2011}, an agent is rational if she responds to what she takes her reasons to be. It is distinct from the structural sense of ETH-25, in which an agent is rational if her choices cohere with her deepest concerns. The two can come apart, and we return to this in Section~\ref{s:amoralist}.

The definition is not trivially satisfied, and it is not stipulated to secure internalism. It is the condition that distinguishes an agent who acts on her view of her reasons from one whose view of her reasons is idle. We argue for its standing as a requirement of rationality in Section~\ref{s:resp}.

\section{The argument}

\begin{theorem}[Internalism for rational agents]\label{t:main}
Necessarily, if a practically rational agent judges that she morally ought to $\phi$, she is motivated to some degree to $\phi$.
\end{theorem}

\begin{proof}
Let the agent judge that she morally ought to $\phi$. By Definition~\ref{d:judgment}, she believes that she morally ought to $\phi$; by \ref{p:reasons}, what she believes entails that there are reasons for her to $\phi$, and, possessing the concept, she believes this. By Definition~\ref{d:rational}, since she is practically rational, she is motivated to some degree to $\phi$. The premises hold necessarily, so the conclusion does.
\end{proof}

\begin{proposition}[Amoralists are possible]\label{p:possible}
There can be an agent who judges that she morally ought to $\phi$ and is not motivated to any degree to $\phi$.
\end{proposition}

\begin{proof}
A moral judgment is a belief, and a motivation is a disposition to act. That an agent believes that $p$ does not entail that she has any particular disposition to act, since the connection between belief and action runs through the agent's practical functioning, which can be impaired while her beliefs remain intact. The cases of acquired sociopathy are actual instances: patients whose capacity to judge what morality requires is unimpaired, as their performance on moral reasoning tasks shows, and whose motivation to act accordingly is absent \citep{Roskies2003}. Nothing in the content of the judgment, as \ref{p:reasons} specifies it, requires the motivation to be present.
\end{proof}

\begin{corollary}[Amoralists are practically irrational]\label{c:irrational}
Every amoralist, an agent who judges that she morally ought to $\phi$ and is not motivated to any degree to $\phi$, is practically irrational.
\end{corollary}

\begin{proof}
By the contrapositive of Theorem~\ref{t:main}: if she were practically rational, she would be motivated to some degree to $\phi$.
\end{proof}

The three results together are the thesis. Theorem~\ref{t:main} is a necessary connection, as internalism requires, but it is conditional on rationality; Proposition~\ref{p:possible} concedes to the externalist that the connection between judgment and motivation is not built into the judgment; and Corollary~\ref{c:irrational} locates the amoralist's defect where the internalist has always felt it to be, in a failure to respond to her own view of what she has reason to do.

\section{Why responsiveness is a requirement of rationality}\label{s:resp}

Definition~\ref{d:rational} carries the argument, and the externalist will say that it is false of rationality: a rational agent, on the Humean view, is one who takes the means to her ends, and an agent without moral ends is not irrational for failing to pursue them. We give three reasons to accept it.

First, the theoretical parallel. An agent who believes that she has conclusive evidence that $p$, and who has no inclination whatever to believe that $p$, is irrational: her doxastic state fails to respond to her own assessment of her evidence. No one takes the Humean line here, holding that she is rational unless she has a desire to believe what her evidence supports. The practical case has the same structure. A belief that there is reason for one to act is an assessment of the considerations bearing on action, and an agent whose motivations do not respond to it at all treats her own assessment as idle.

Second, the explanatory fact that \citet{Smith1994} called striking: in a good and strong-willed person, a change in moral judgment is reliably followed by a change in motivation. If an agent comes to believe that she ought to give to famine relief, she becomes motivated to give; if she is persuaded that she was wrong, the motivation lapses. The externalist explains the regularity by a standing desire to do what is right, and Smith objected that this makes the good person a fetishist, moved by rightness as such rather than by the features that make acts right. Definition~\ref{d:rational} explains the regularity without the fetish. The rational agent is moved by the considerations she takes to be reasons, by the hunger of the famine victims and the difference her gift would make, and her motivation changes when her view of these considerations changes.

Third, the account of reasons. MTE-11 argues that reasons are primitive in the normative order, and ACT-12 that there are external reasons, which an agent may have although no sound deliberation from her motivations would reveal them. On the Humean view, an agent's reasons are fixed by her motivations, and the requirement of rationality is to act on them. If there are external reasons, the Humean picture of rationality loses its rationale: rationality cannot consist in serving one's motivations if reasons are not grounded in them. What remains is the requirement that the agent's motivations respond to what she takes her reasons to be, which is Definition~\ref{d:rational}.

\section{The amoralist}\label{s:amoralist}

Corollary~\ref{c:irrational} classifies, and the cases differ.

\emph{The inverted-commas amoralist} says that an act is right, meaning that it is what others call right \citep{Hare1952}. He does not make the moral judgment, and he is not a counterexample to anything. He is not irrational either, since he has no belief about his reasons that his motivations fail to answer to.

\emph{The depressed agent} judges that she ought to visit her sick friend, and cannot bring herself to care \citep{Stocker1979}. She makes the judgment and lacks the motivation, and her condition is a paradigm of impaired practical functioning. Corollary~\ref{c:irrational} classifies her as practically irrational, and the classification is correct in the present sense, though it carries no blame: practical irrationality is a defect in the relation between an agent's normative beliefs and her motivations, and some such defects are illnesses.

\emph{The patient with acquired sociopathy} makes moral judgments with normal competence after a brain injury and is no longer moved by them \citep{Roskies2003}. This is the actual amoralist, and Proposition~\ref{p:possible} uses him. He is irrational in the sense of Definition~\ref{d:rational}, since he believes that there are reasons for him to act and does not respond to them, although his means-end reasoning is intact.

\emph{The philosophical amoralist} of \citet{Brink1989} understands morality, makes moral judgments and denies that they give him any reason. If he believes that he morally ought to $\phi$ and that he has no reason to $\phi$, he has inconsistent beliefs by \ref{p:reasons}, and he is theoretically irrational. If he makes the moral judgment without the second belief, he is practically irrational by Corollary~\ref{c:irrational}. Either way he is irrational, though not, perhaps, in the same way.

The amoralist may be rational in the structural sense of ETH-25. An agent without any concern for others or for morality, whose choices cohere with the concerns she has, makes no mistake by the standard of coherence with her deepest concerns. The thesis does not deny this; it says that she makes a mistake by the standard of responsiveness to her own view of her reasons. The two standards coincide in rational agents who accept a moral requirement, since such an agent is moved by the requirement and so has a concern for meeting it, which is the consequence ETH-25 draws.

\section{Objections and replies}

\begin{objection}[Triviality]
The thesis is trivial. If being practically rational is defined as being motivated by one's normative judgments, then it is analytic that rational agents are motivated by them, and the dispute with the externalist has been settled by stipulation.
\end{objection}

\begin{reply}
Definition~\ref{d:rational} does not mention moral judgment. It concerns beliefs about reasons in general, and it is supported by considerations, the theoretical parallel and the striking fact, that are independent of the dispute about moral motivation. The internalist conclusion follows only with \ref{p:reasons}, which is a substantive thesis about the content of moral judgments and is denied by those who hold that moral requirements need not give reasons \citep{Foot1972}. The argument is therefore not trivial: it derives the internalist connection from an account of rationality and an account of moral ought, each defended on its own grounds.
\end{reply}

\begin{objection}[Psychopaths are not irrational]
Psychopaths reason well about means and ends, and many are highly intelligent. To call them irrational is to use the word for whatever one disapproves of \citep{Nichols2002}.
\end{objection}

\begin{reply}
Psychopaths are not irrational in the sense of means-end reasoning, and the thesis does not say they are. It says that, if they make moral judgments, they are irrational in the sense of Definition~\ref{d:rational}: their motivations fail to respond to what they believe to be reasons for them. This is a determinate condition, distinct from disapproval. Whether psychopaths do make moral judgments is an empirical question on which the evidence is mixed; if they do not, because they lack full competence with moral concepts, they are not amoralists, and the thesis does not classify them as irrational at all.
\end{reply}

\begin{objection}[The Humean]
Reasons for action are given by an agent's desires, and an agent who lacks moral desires lacks moral reasons. The amoralist is therefore not unmoved by reasons she has; she simply has none, and her moral belief, if she is a Humean, is false.
\end{objection}

\begin{reply}
The objection denies \ref{p:reasons} in its second clause, that moral reasons are reasons for every agent to whom they apply. ACT-12 argues that external reasons exist, and MTE-11 that the reason relation is identical to a natural relation specified by idealized reactions that do not depend on the agent's own motivations. If these arguments succeed, the amoralist has the moral reasons her judgment represents. If they fail, the Humean amoralist who believes that she morally ought to $\phi$ believes that there is a decisive consideration bearing on her action that gives her no reason, and the question would be what the moral ought then means. Either way, the objection does not reinstate the amoralist as a rational agent who makes the moral judgment in the sense of \ref{p:reasons}.
\end{reply}

\begin{objection}[Weakness of will]
Akratic agents judge that they ought to act and do not act. If rational agents are moved by their moral judgments, akrasia should be impossible for them; but akrasia is common among agents who are otherwise rational.
\end{objection}

\begin{reply}
Theorem~\ref{t:main} requires only motivation to some degree, and the akratic agent has it: she is moved by her judgment, and more strongly moved by something else. Her failure is a failure of the strength of her motivation relative to her judgment of the weight of her reasons. On a natural strengthening of Definition~\ref{d:rational}, on which a rational agent's motivations are proportioned to her assessment of the strength of her reasons, the akratic agent is also practically irrational, as is usually thought. The thesis requires only the weak form, and is consistent with ACT-04.
\end{reply}

\section{What the result does not establish}

The thesis does not establish that moral judgments are beliefs; it assumes it, with MTE-01, and it would need restatement on an expressivist view. It does not establish that amoralists are blameworthy for their condition, since practical irrationality may be an illness or an injury. It does not say how strongly rational agents are moved by their moral judgments, only that they are moved to some degree, and so it does not entail that rational agents always do what they judge they ought to do. Nor does it say that rational agents have moral motivation independently of their judgments: an agent who has no moral beliefs is not irrational for lacking moral motivation, though, if MTE-08 is right that moral truths are known by a generally reliable faculty, she may be failing to form beliefs she has every opportunity to form.

\section{Conclusion}

Moral judgment necessarily motivates rational agents. To judge that one morally ought to act is to judge that there are decisive moral reasons for one to act, and a practically rational agent is moved, to some degree, by what she takes to be reasons for her. The connection holds of necessity, but it runs through rationality rather than through the content of the judgment: the judgment is a belief, and an agent can have it without the motivation. That agent is the amoralist. She is possible, and some actual agents after brain injury are instances; and she is practically irrational, because her motivations fail to answer to her own view of her reasons.

\begin{thebibliography}{99}\small
\bibitem[Brink(1989)]{Brink1989} Brink, D.~O. (1989). \emph{Moral Realism and the Foundations of Ethics}. Cambridge University Press.
\bibitem[Foot(1972)]{Foot1972} Foot, P. (1972). Morality as a system of hypothetical imperatives. \emph{The Philosophical Review}, 81(3), 305--316.
\bibitem[Frankena(1958)]{Frankena1958} Frankena, W.~K. (1958). Obligation and motivation in recent moral philosophy. In A.~I. Melden (Ed.), \emph{Essays in Moral Philosophy}. University of Washington Press.
\bibitem[Hare(1952)]{Hare1952} Hare, R.~M. (1952). \emph{The Language of Morals}. Oxford University Press.
\bibitem[Hume(1739)]{HumeTreatise} Hume, D. (1739--1740). \emph{A Treatise of Human Nature}.
\bibitem[Korsgaard(1986)]{Korsgaard1986} Korsgaard, C.~M. (1986). Skepticism about practical reason. \emph{The Journal of Philosophy}, 83(1), 5--25.
\bibitem[Nichols(2002)]{Nichols2002} Nichols, S. (2002). How psychopaths threaten moral rationalism: Is it irrational to be amoral? \emph{The Monist}, 85(2), 285--303.
\bibitem[Parfit(2011)]{Parfit2011} Parfit, D. (2011). \emph{On What Matters}, vol.~1. Oxford University Press.
\bibitem[Roskies(2003)]{Roskies2003} Roskies, A. (2003). Are ethical judgments intrinsically motivational? Lessons from \enquote{acquired sociopathy}. \emph{Philosophical Psychology}, 16(1), 51--66.
\bibitem[Smith(1994)]{Smith1994} Smith, M. (1994). \emph{The Moral Problem}. Blackwell.
\bibitem[Stocker(1979)]{Stocker1979} Stocker, M. (1979). Desiring the bad: An essay in moral psychology. \emph{The Journal of Philosophy}, 76(12), 738--753.
\bibitem[Svavarsd{\'o}ttir(1999)]{Svavarsdottir1999} Svavarsd{\'o}ttir, S. (1999). Moral cognitivism and motivation. \emph{The Philosophical Review}, 108(2), 161--219.
\end{thebibliography}

\end{document}
